\clearpage
\phantomsection
\section*{Chương 0. Bài toán, mục tiêu so sánh, và Praat thật sự làm gì}
\addcontentsline{toc}{section}{Chương 0. Bài toán, mục tiêu so sánh, và Praat thật sự làm gì}

\begin{learningobj}
\begin{itemize}[leftmargin=1.5em, itemsep=2pt]
    \item Nắm vững bản chất bài toán ước lượng tần số cơ bản ($F_0$) và lý do vì sao hai phương pháp trực giác (đếm điểm qua không và tìm đỉnh phổ FFT) thất bại trên tín hiệu tiếng nói thực tế.
    \item Hiểu rõ nguyên lý cốt lõi của Boersma: chuẩn hóa hàm tự tương quan bằng cách chia cho hàm tự tương quan của cửa sổ Hanning để triệt tiêu độ suy giảm giả tạo ở các độ trễ lớn.
    \item Nắm tường tận 8 bước thuật toán trong mã nguồn Praat 7.0.03 cùng 10 tham số chuẩn, đối chiếu chính xác giữa các khâu phần mềm của Praat và các khối xử lý phần cứng trên FPGA.
    \item Hiểu rõ kiến trúc toàn hệ thống: luồng dữ liệu từ tệp WAV qua HPS, cầu nối Lightweight HPS-to-FPGA, lõi tính toán trên FPGA, và đường trả kết quả F0 về HPS để ghi tệp nhị phân.
    \item Nắm chắc định nghĩa 4 chỉ số sai số chuẩn trong so sánh cao độ: VDE, GPE, FFE và sai số cent mịn, cùng cách tính tay trên từng khung.
    \item Phân tích ngân sách chu kỳ từ nguyên lý đầu tiên: 50 MHz, 3.125 chu kỳ/mẫu, 500.000 chu kỳ/bước nhảy 10 ms; dẫn xuất chính xác 98.540 phép nhân-cộng MAC và đo thực tế 98.543 chu kỳ/khung, từ đó xác lập hệ số thời gian thực (RTF).
    \item Hiểu kiến trúc kiểm chứng ba tầng và vai trò của bộ mô hình tương đương (twin) số nguyên \texttt{golden\_v2/fixed\_twin.py} làm đặc tả phần cứng cho toàn bộ dự án.
\end{itemize}
\end{learningobj}

\subsection*{0.1. Một file, hai câu trả lời}
\addcontentsline{toc}{subsection}{0.1. Một file, hai câu trả lời}

Một tệp âm thanh WAV chuẩn 16 kHz, kênh đơn (mono), định dạng mẫu PCM 16-bit có dấu được đưa vào hệ thống. Trên máy tính cá nhân, phần mềm Praat đọc tệp âm thanh này và xuất ra một đường cao độ $F_0$ với chu kỳ phân tích 10 ms. Trên bo mạch DE10-Standard, hệ thống vi xử lý HPS (Hard Processor System, Cortex-A9 lõi kép) đọc cùng tệp WAV đó từ thẻ nhớ SD, truyền các mẫu âm thanh qua cầu nối bus sang vùng logic FPGA (Cyclone V 5CSXFC6D6F31C6N). Lõi phần cứng trên FPGA tính toán cao độ từng khung và trả kết quả về HPS dưới dạng tệp nhị phân \texttt{pitch\_f0.bin} ở định dạng số cố định Q12.4.

Câu hỏi trung tâm mà anh Tú cùng bất kỳ kỹ sư phần cứng nào cũng phải đối mặt là: hai đường cong cao độ đó lệch nhau bao nhiêu, và nguyên nhân vật lý hay số học nào tạo ra sự chênh lệch ở từng mili-giây? Để trả lời thấu đáo câu hỏi này, trước hết ta phải hiểu vì sao việc ước lượng cao độ tiếng nói lại khó, và vì sao những cách tiếp cận trực giác thông thường đều sụp đổ khi áp dụng vào thực tế.

\subsubsection*{Sự sụp đổ của phương pháp đếm điểm qua không (Zero-Crossing Counting)}
Phương pháp đơn giản nhất mà người mới học xử lý tín hiệu thường nghĩ đến là đếm số lần tín hiệu đổi dấu qua mức 0 (zero-crossing) trong một khoảng thời gian, từ đó suy ra tần số. Tuy nhiên, một sự thật vật lý cơ bản là: tiếng nói có thanh (voiced speech) của con người hầu như không bao giờ là một sóng hình sin đơn âm. Khi dây thanh âm rung động, luồng khí đóng mở theo chu kỳ tạo ra dạng sóng có chứa rất nhiều hài âm (harmonics). Nếu hài âm bậc hai ($2 f_0$) có biên độ mạnh và lệch pha so với sóng cơ bản, nó sẽ uốn dạng sóng tổng hợp cắt qua trục 0 nhiều lần trong cùng một chu kỳ lặp lại.

Để thấy rõ sự thất bại này bằng số liệu định lượng, ta xét một đoạn tín hiệu tổng hợp được tạo bởi kịch bản kiểm chứng \texttt{verify\_v3/ch00\_naive\_methods.py}:
\begin{equation}
x(t) = \sin(2\pi \cdot 100 t) - 0{,}75 \cos(2\pi \cdot 200 t)
\end{equation}
Tần số cơ bản thực sự là $f_0 = 100{,}0\text{ Hz}$, tần số lấy mẫu $f_s = 16.000\text{ Hz}$. Một chu kỳ cơ bản kéo dài $T_0 = 1/100 = 10\text{ ms}$, tương ứng đúng $P = 160$ mẫu. Nếu tín hiệu là sóng sin thuần, một chu kỳ 160 mẫu sẽ cắt qua trục 0 đúng 2 lần. Nhưng khi kiểm tra trên đoạn tín hiệu dài 3 chu kỳ (480 mẫu), bộ đếm điểm qua không ghi nhận tổng cộng 12 lần đổi dấu, tức trung bình 4 điểm qua không trong mỗi chu kỳ. Tần số ước lượng bị thổi phồng thành:
\begin{equation}
f_{\text{est}} = \frac{12}{2 \times 3 \times 0{,}01\text{ s}} = 200{,}0\text{ Hz}
\end{equation}
Thuật toán bị đánh lừa hoàn toàn bởi hài âm bậc hai, báo sai lệch cao độ gấp đôi ($100\%$ lỗi) ngay trên một tín hiệu hoàn toàn tuần hoàn và không hề có nhiễu bên ngoài.

\textbf{Ví dụ tính tay thu nhỏ.} Để nhìn thấy từng bước ở mức mẫu, xét 16 mẫu rời rạc đại diện cho một chu kỳ rút gọn (lấy trực tiếp từ log của \texttt{verify\_v3/ch00\_naive\_methods.py}):
\begin{center}
\small
\begin{tabular}{lrrrrrrrrrrrrrrrr}
\toprule
$n$ & 0 & 1 & 2 & 3 & 4 & 5 & 6 & 7 & 8 & 9 & 10 & 11 & 12 & 13 & 14 & 15 \\
\midrule
$x[n]$ & 0 & $-18$ & $-9$ & 36 & 100 & 149 & 151 & 95 & 0 & $-95$ & $-151$ & $-149$ & $-100$ & $-36$ & 9 & 18 \\
\bottomrule
\end{tabular}
\end{center}
Quan sát dãy mẫu: từ $n=2$ (giá trị $-9$) sang $n=3$ ($+36$), tín hiệu đổi dấu lần thứ nhất. Đến $n=8$, tín hiệu bằng 0 và giảm xuống $-95$ ở $n=9$, đổi dấu lần thứ hai. Từ $n=13$ ($-36$) sang $n=14$ ($+9$), tín hiệu lại đổi dấu lần thứ ba. Cùng với điểm 0 ban đầu ở $n=0$, ta có 4 lần cắt qua 0 trong chỉ một chu kỳ cơ bản. Mạch đếm phần cứng nếu dựa vào hiện tượng này sẽ nhận diện nhầm chu kỳ thành 8 mẫu thay vì 16 mẫu.

\subsubsection*{Sự sụp đổ của phương pháp chọn đỉnh phổ FFT (FFT Peak Picking)}
Cách tiếp cận thứ hai là chuyển tín hiệu sang miền tần số bằng phép biến đổi Fourier nhanh (FFT, Fast Fourier Transform) rồi tìm vị trí bin có biên độ cực đại. Để bám sát ngữ điệu biến đổi liên tục của tiếng nói, ta bắt buộc phải phân tích trên từng khung thời gian ngắn khoảng vài chục mili-giây (thường từ 20 ms đến 40 ms). Tại tần số lấy mẫu $f_s = 16.000\text{ Hz}$, một khung 40 ms tương ứng với 640 mẫu (hoặc 638 mẫu theo quy ước làm tròn đối xứng của Praat).

Độ phân giải tần số của phép phân tích FFT bị chặn cứng bởi chiều dài khung $N$:
\begin{equation}
\Delta f = \frac{f_s}{N} = \frac{16.000}{640} = 25{,}00\text{ Hz} \quad (\text{với } N=638, \Delta f = \frac{16.000}{638} \approx 25{,}08\text{ Hz})
\end{equation}
Khoảng cách giữa hai vạch phổ kề nhau lên tới 25 Hz. Hãy xét một giọng nam trầm có tần số cơ bản thực là $f_0 = 112{,}5\text{ Hz}$. Tần số này rơi vào chính giữa bin 4 ($4 \times 25 = 100\text{ Hz}$) và bin 5 ($5 \times 25 = 125\text{ Hz}$). Bộ chọn đỉnh phổ dù hoàn hảo nhất cũng chỉ có thể chỉ định bin 4 hoặc bin 5, gây ra sai số tuyệt đối ít nhất là $12{,}5\text{ Hz}$.

Trong âm nhạc học và ngữ âm học, tai người cảm nhận cao độ theo hàm logarit của tần số. Đơn vị đo độ chênh lệch cao độ chuẩn là \textbf{cent}, trong đó 1 quãng tám (tần số gấp đôi) tương ứng với 1.200 cent, và 1 nửa cung (semitone) tương ứng với 100 cent:
\begin{equation}
\Delta \text{cent} = 1.200 \cdot \left| \log_2 \left( \frac{f_{\text{est}}}{f_0} \right) \right|
\end{equation}
Áp dụng cho trường hợp chọn nhầm bin 100 Hz khi tần số thực là 112,5 Hz:
\begin{equation}
\Delta \text{cent} = 1.200 \cdot \left| \log_2 \left( \frac{100{,}0}{112{,}5} \right) \right| \approx 203{,}91\text{ cent}
\end{equation}
(kết quả được kiểm chứng trong \texttt{verify\_v3/ch00\_naive\_methods.log}). Sai số hơn 200 cent tương đương với hai nửa cung (một cung nguyên) trong âm nhạc. Một bộ ước lượng $F_0$ với sai số thô như vậy hoàn toàn vô giá trị trong các ứng dụng phân tích ngữ điệu hoặc đối sánh khoa học. Muốn giảm $\Delta f$ xuống dưới 1 Hz, FFT cần cửa sổ dài hơn 1 giây, nhưng trong 1 giây đó người nói đã phát âm qua 3 đến 5 âm tiết khác nhau, tính dừng của tín hiệu bị phá vỡ hoàn toàn.

\subsubsection*{Bản chất của phép tự tương quan (Autocorrelation)}
Để khắc phục triệt để hai nhược điểm trên, ta quay về nguyên lý đầu tiên: \textbf{một tín hiệu tuần hoàn luôn có hình dạng giống hệt chính nó khi bị dịch trễ đi đúng một chu kỳ cơ bản}. Phép đo mức độ tương đồng giữa tín hiệu $x[n]$ và bản sao dịch trễ của nó $x[n+\tau]$ ở độ trễ $\tau$ (lag) được định nghĩa qua hàm tự tương quan ngắn hạn:
\begin{equation}
R[\tau] = \sum_{n=0}^{N-1-\tau} x[n] x[n+\tau]
\end{equation}

\textbf{Ví dụ tính tay số nguyên thu nhỏ.} Xét một tín hiệu rời rạc tuần hoàn hoàn hảo với chu kỳ cơ bản $P = 3$ mẫu, độ dài khung $N = 9$ mẫu:
\begin{equation}
x = [3, -1, -2, \; 3, -1, -2, \; 3, -1, -2]
\end{equation}
Ta tính trực tiếp giá trị tích vô hướng $R[\tau]$ theo từng bước cho các độ trễ từ $\tau = 0$ đến $\tau = 5$ (trích xuất từ kịch bản kiểm chứng \texttt{verify\_v3/ch00\_naive\_methods.py}):
\begin{itemize}[leftmargin=1.5em, itemsep=2pt]
    \item Tại $\tau = 0$: lấy 9 cặp mẫu tự nhân với chính nó:
    \begin{align*}
    R[0] &= 3^2 + (-1)^2 + (-2)^2 + 3^2 + (-1)^2 + (-2)^2 + 3^2 + (-1)^2 + (-2)^2 \\
         &= 9 + 1 + 4 + 9 + 1 + 4 + 9 + 1 + 4 = 42
    \end{align*}
    \item Tại $\tau = 1$: nhân 8 cặp mẫu kề nhau:
    \begin{align*}
    R[1] &= 3(-1) + (-1)(-2) + (-2)3 + 3(-1) + (-1)(-2) + (-2)3 + 3(-1) + (-1)(-2) \\
         &= -3 + 2 - 6 - 3 + 2 - 6 - 3 + 2 = -15
    \end{align*}
    \item Tại $\tau = 2$: nhân 7 cặp mẫu cách nhau 2 vị trí:
    \begin{align*}
    R[2] &= 3(-2) + (-1)3 + (-2)(-1) + 3(-2) + (-1)3 + (-2)(-1) + 3(-2) \\
         &= -6 - 3 + 2 - 6 - 3 + 2 - 6 = -20
    \end{align*}
    \item Tại $\tau = 3$: nhân 6 cặp mẫu cách nhau đúng một chu kỳ $P=3$:
    \begin{align*}
    R[3] &= 3(3) + (-1)(-1) + (-2)(-2) + 3(3) + (-1)(-1) + (-2)(-2) \\
         &= 9 + 1 + 4 + 9 + 1 + 4 = 28
    \end{align*}
    \item Tại $\tau = 4$: nhân 5 cặp mẫu, thu được $R[4] = -8$.
    \item Tại $\tau = 5$: nhân 4 cặp mẫu, thu được $R[5] = -13$.
\end{itemize}

\textbf{Chỗ dễ vướng: Vì sao $R[3] = 28$ lại bé hơn $R[0] = 42$ dù tín hiệu tuần hoàn tuyệt đối?}
Đây là bẫy tư duy kinh điển mà người thiết kế phần cứng hay gặp. Ở độ trễ $\tau = 3$, số cặp mẫu còn lại để nhân với nhau chỉ là $N - \tau = 9 - 3 = 6$ cặp, trong khi ở $\tau = 0$ ta cộng dồn tới 9 cặp. Nếu ta tính giá trị trung bình không chệch bằng cách chia tổng tích lũy cho số cặp mẫu tham gia:
\begin{align}
\frac{R[0]}{9 - 0} &= \frac{42}{9} \approx 4{,}67 \\
\frac{R[3]}{9 - 3} &= \frac{28}{6} \approx 4{,}67
\end{align}
Hai con số hoàn toàn trùng khớp! Đỉnh tương quan tại $\tau = 3$ đã chỉ ra chính xác chu kỳ cơ bản của tín hiệu là 3 mẫu mà không bị nhiễu hài âm đánh lừa như phương pháp đếm điểm qua 0, cũng không bị lượng tử hóa thô thiển như phương pháp chọn đỉnh FFT.

\textit{Tóm lại: Đếm điểm qua không thất bại vì hài âm bậc hai tạo thêm điểm cắt giả, lấy đỉnh FFT thất bại vì độ phân giải bin thô tới 204 cent ở 100 Hz; phương pháp tự tương quan đo mức độ tương đồng theo miền thời gian, mở ra con đường xác định chu kỳ cơ bản chính xác và bền vững trước cấu trúc hài âm phức tạp của tiếng nói.}

\subsection*{0.2. Ý tưởng của Boersma: Chuẩn hóa bằng tự tương quan của cửa sổ}
\addcontentsline{toc}{subsection}{0.2. Ý tưởng của Boersma: Chuẩn hóa bằng tự tương quan của cửa sổ}

Khi áp dụng phương pháp tự tương quan lên các đoạn tiếng nói thực tế, một bài toán kỹ thuật mới lập tức nảy sinh: ta không thể tính toán trên toàn bộ tệp âm thanh dài hàng phút mà bắt buộc phải cắt tín hiệu thành từng khung ngắn hữu hạn (khoảng 40 ms) để theo dõi sự biến đổi cao độ theo thời gian. Việc cắt khung bằng một hàm chữ nhật đột ngột (rectangular window) sẽ tạo ra các điểm gián đoạn sắc nhọn ở hai biên khung, dẫn đến hiện tượng rò rỉ năng lượng phổ (spectral leakage) nghiêm trọng. Để khắc phục hiện tượng này, giải pháp kinh điển là nhân khung tín hiệu với một hàm cửa sổ êm dịu, tiêu biểu là cửa sổ Hanning (hoặc Hamming), có biên độ vuốt dần về 0 ở hai đầu mút.

Tuy nhiên, việc nhân cửa sổ lại kéo theo một hậu quả tiêu cực khác: các mẫu âm thanh nằm xa tâm khung bị suy giảm biên độ mạnh mẽ. Khi ta tính hàm tự tương quan trên tín hiệu đã nhân cửa sổ $x_w[n] = x[n] \cdot w[n]$:
\begin{equation}
R_{xw}[\tau] = \sum_{n=0}^{N-1-\tau} x_w[n] x_w[n+\tau] = \sum_{n=0}^{N-1-\tau} x[n] w[n] x[n+\tau] w[n+\tau]
\end{equation}
Độ trễ $\tau$ càng lớn thì khoảng chồng lấn giữa $w[n]$ và $w[n+\tau]$ càng thu hẹp, và tích $w[n] w[n+\tau]$ càng bị đè bẹp về gần 0. Đồ thị $R_{xw}[\tau]$ xuất hiện một sườn dốc suy thoái nhân tạo theo $\tau$. Hệ quả là đỉnh tương quan ứng với chu kỳ cơ bản thực sự $T_0$ của một giọng nam trầm (nằm ở các độ trễ lớn, ví dụ $\tau \approx 160$ mẫu tại 100 Hz) bị lùn đi nghiêm trọng. Nếu dùng hàm tương quan chuẩn hóa thông thường $R_{xw}[\tau]/R_{xw}[0]$, giá trị đỉnh này thường xuyên rơi xuống dưới ngưỡng hữu thanh (voicing threshold mặc định của Praat là 0,45) khiến khung âm thanh bị nhận diện nhầm thành vô thanh, hoặc bị lép vế trước một đỉnh phụ ở độ trễ nhỏ hơn ($\tau = T_0 / 2$), gây ra lỗi kinh điển là \textbf{nhân đôi quãng tám (octave doubling)}.

\subsubsection*{Cơ chế chia cho tự tương quan cửa sổ $R_w[\tau]$}
Đóng góp đột phá của Paul Boersma (1993) nằm ở nhận định giải tích: hàm tự tương quan của tín hiệu đã nhân cửa sổ xấp xỉ bằng tích giữa hàm tự tương quan thực sự của tín hiệu $R_x[\tau]$ với hàm tự tương quan của chính hàm cửa sổ $R_w[\tau]$:
\begin{equation}
R_{xw}[\tau] \approx R_x[\tau] \cdot R_w[\tau]
\end{equation}
Trong đó hàm tự tương quan của cửa sổ được định nghĩa hoàn toàn tương tự:
\begin{equation}
R_w[\tau] = \sum_{n=0}^{N-1-\tau} w[n] w[n+\tau]
\end{equation}
Từ mối liên hệ nhân này, Boersma đề xuất thuật toán chuẩn hóa: muốn tìm lại hàm tự tương quan chuẩn hóa không thiên vị $r[\tau]$ của bản thân tín hiệu gốc (với giá trị lý tưởng $r[0] = 1$ và $r[T_0] = 1$ đối với tín hiệu tuần hoàn hoàn hảo), ta chỉ cần lấy $R_{xw}[\tau]$ chia cho $R_w[\tau]$, rồi chuẩn hóa tỉ lệ theo năng lượng tại $\tau = 0$:
\begin{equation}
r[\tau] = \frac{R_{xw}[\tau] / R_w[\tau]}{R_{xw}[0] / R_w[0]} = \frac{R_{xw}[\tau]}{R_w[\tau]} \cdot \frac{R_w[0]}{R_{xw}[0]}
\end{equation}
Phép chia này đóng vai trò như một bộ bù độ dốc: vì $R_w[\tau]$ giảm dần về 0 khi $\tau$ tăng theo đúng quy luật hình học của cửa sổ, việc chia cho $R_w[\tau]$ sẽ nâng bổng các đỉnh ở độ trễ xa lên đúng biên độ thực sự của chúng, triệt tiêu hoàn toàn sự thiên lệch về phía các độ trễ nhỏ.

\subsubsection*{Vì sao phép xấp xỉ $R_{xw}[\tau] \approx R_x[\tau] \cdot R_w[\tau]$ lại đúng?}
Để trả lời thuyết phục người kỹ sư phần cứng, ta xét bản chất vật lý của hai chuyển động: dao động nhanh của dây thanh âm và sự biến thiên chậm của đường bao cửa sổ. 

Trong thuật toán Praat, chiều dài khung phân tích được ấn định bằng đúng 3 chu kỳ của tần số thấp nhất (pitch floor $f_{\min} = 75\text{ Hz}$):
\begin{equation}
T_{\text{window}} = \frac{3}{f_{\min}} = \frac{3}{75} = 0{,}04\text{ s} = 40\text{ ms} \quad (N = 638\text{ mẫu tại } f_s = 16\text{ kHz})
\end{equation}
Trong khi đó, tần số cơ bản của người nói dao động từ 75 Hz đến 600 Hz, tương ứng với chu kỳ cơ bản $T_0$ từ 26 mẫu (600 Hz) đến 213 mẫu (75 Hz). Điều này có nghĩa là bên trong một khung 638 mẫu, tín hiệu tiếng nói thực hiện từ 3 đến hơn 24 chu kỳ dao động tuần hoàn nhanh.

Đường bao cửa sổ Hanning $w[n]$ là một hàm cosin trơn dốc êm dịu trải dài trên toàn bộ 638 mẫu. Trong phạm vi một chu kỳ đơn lẻ $T_0$, giá trị của $w[n]$ và $w[n+\tau]$ gần như là một hằng số cục bộ. Khi ta lấy tổng tích lũy $\sum x[n] x[n+\tau] \cdot (w[n] w[n+\tau])$, đại lượng dao động nhanh $x[n] x[n+\tau]$ tạo ra giá trị kỳ vọng trung bình tuần hoàn $R_x[\tau]$, trong khi đại lượng biến thiên chậm $w[n] w[n+\tau]$ có thể đưa ra ngoài dấu tích phân cục bộ dưới dạng trọng số bao $R_w[\tau]$. Phép nhân chập giữa phổ tín hiệu hẹp và phổ cửa sổ rộng trong miền tần số tương ứng với phép nhân đường bao trực tiếp trong miền thời gian.

\subsubsection*{Ví dụ tính tay thu nhỏ: Sóng sin thuần qua cửa sổ Hanning 8 mẫu}
Để minh chứng tường minh bằng số học cụ thể, ta xét một thí nghiệm mô phỏng thu nhỏ được thực thi bởi kịch bản \texttt{verify\_v3/ch00\_norm\_sine.py}. Ta sử dụng một sóng sin tuần hoàn hoàn hảo có chu kỳ $P = 4$ mẫu (tương ứng tần số $f_0 = f_s / 4 = 4.000\text{ Hz}$), chiều dài phân tích $N = 8$ mẫu:
\begin{equation}
x[n] = \sin\left(\frac{2\pi n}{4}\right) = [0{,}000, \; 1{,}000, \; 0{,}000, \; -1{,}000, \; 0{,}000, \; 1{,}000, \; 0{,}000, \; -1{,}000]
\end{equation}
Hàm cửa sổ Hanning 8 mẫu của Praat tuân thủ công thức $N+1$ (được định nghĩa trong mã nguồn Praat \texttt{Sound\_to\_Pitch.cpp} dòng 423--424):
\begin{equation}
w[i] = 0{,}5 - 0{,}5 \cos\left(\frac{2\pi i}{N+1}\right), \quad i = 1, \dots, N \quad (i = n + 1)
\end{equation}
Với $N=8$, các mẫu cửa sổ tính ra là:
\begin{equation*}
w = [0{,}116978, \; 0{,}413176, \; 0{,}750000, \; 0{,}969846, \; 0{,}969846, \; 0{,}750000, \; 0{,}413176, \; 0{,}116978]
\end{equation*}

Tín hiệu sau khi nhân cửa sổ $x_w[n] = x[n] \cdot w[n]$ là:
\begin{align*}
x_w = [& 0{,}0000, \; 0{,}4132, \; 0{,}0000, \; -0{,}9698, \; -0{,}0000, \; 0{,}7500, \; 0{,}0000, \; -0{,}1170]
\end{align*}
Bảng \ref{tab:ch00_norm_sine} liệt kê kết quả tính tay chi tiết hàm tự tương quan của cửa sổ $R_w[\tau]$, hàm tự tương quan của tín hiệu nhân cửa sổ $R_{xw}[\tau]$, hệ số tương quan chưa chuẩn hóa $r_{\text{unnorm}}[\tau] = R_{xw}[\tau]/R_{xw}[0]$, và hệ số chuẩn hóa Boersma $r_{\text{Boersma}}[\tau]$ cho các độ trễ từ $\tau = 0$ đến $\tau = 4$ (trích xuất nguyên văn từ \texttt{verify\_v3/ch00\_norm\_sine.log}).

\begin{table}[H]
\centering
\small
\caption{Đối chiếu giá trị tự tương quan trước và sau khi chuẩn hóa cửa sổ Boersma trên sóng sin thuần ($N=8, P=4$)}
\label{tab:ch00_norm_sine}
\begin{tabular}{crrrrc}
\toprule
Độ trễ $\tau$ & $R_w[\tau]$ & $R_{xw}[\tau]$ & $r_{\text{unnorm}}[\tau]$ & $r_{\text{Boersma}}[\tau]$ & Đánh giá vật lý \\
\midrule
0 & 3,375000 & 1,687500 & 1,000000 & 1,000000 & Gốc chuẩn hóa năng lượng \\
1 & 3,111800 & 0,000000 & 0,000000 & 0,000000 & Trực giao ở góc pha $90^\circ$ \\
2 & 2,431670 & $-1{,}215835$ & $-0{,}720495$ & $-1{,}000000$ & Ngược pha hoàn toàn ($180^\circ$) \\
3 & 1,590835 & 0,000000 & 0,000000 & 0,000000 & Trực giao ở góc pha $270^\circ$ \\
\textbf{4} & \textbf{0,846665} & \textbf{0,423332} & \textbf{0,250864} & \textbf{1,000000} & \textbf{Chu kỳ cơ bản $P=4$} \\
\bottomrule
\end{tabular}
\end{table}

\textbf{Đọc bảng và kết luận định lượng.} Hãy nhìn thẳng vào dòng $\tau = 4$, chính là chu kỳ lặp lại thực tế của sóng sin:
\begin{itemize}[leftmargin=1.5em, itemsep=2pt]
    \item Nếu \textbf{không chuẩn hóa cửa sổ}, hệ số tương quan chỉ đạt $r_{\text{unnorm}}[4] = 0{,}250864 \approx 0{,}251$. Tín hiệu thuần nhất 100\% tuần hoàn nhưng giá trị đỉnh tương quan bị sụt giảm tới gần 75\% năng lượng! Với ngưỡng voicing mặc định của Praat là 0,45, thuật toán sẽ kết luận ngay khung này là \emph{vô thanh} (unvoiced) do $0{,}251 < 0{,}45$.
    \item Ngược lại, khi áp dụng \textbf{chuẩn hóa Boersma}: ta lấy $R_{xw}[4] / R_w[4] = 0{,}423332 / 0{,}846665 = 0{,}500000$. Chia cho tỉ số tại gốc $R_{xw}[0] / R_w[0] = 1{,}687500 / 3{,}375000 = 0{,}500000$, ta thu được kết quả chính xác tuyệt đối:
    \begin{equation}
    r_{\text{Boersma}}[4] = \frac{0{,}500000}{0{,}500000} = 1{,}000000
    \end{equation}
\end{itemize}
Giá trị 1,000000 phản ánh đúng 100\% tính tuần hoàn của sóng sin, không hề bị ảnh hưởng bởi việc cửa sổ bị thót hẹp ở hai đầu. Đây chính là linh hồn toán học giải thích vì sao phần mềm Praat đạt được độ tin cậy và danh tiếng vượt trội trong cộng đồng nghiên cứu ngữ âm học suốt ba thập kỷ qua.

\textit{Tóm lại: Cửa sổ cắt khung làm suy hao giả tạo năng lượng ở các độ trễ lớn khiến đỉnh cao độ thật bị lùn đi; ý tưởng cốt lõi của Boersma là chia hàm tự tương quan của tín hiệu cho hàm tự tương quan của cửa sổ $R_{xw}[\tau] / R_w[\tau]$, khôi phục trọn vẹn độ mạnh của chu kỳ cơ bản và triệt tiêu định kiến thiên lệch về các độ trễ nhỏ.}
